\section{Outer routing and T2-C details}
\label{app:routing}

\paragraph{Parameter-only regime routing.}
The E2 router receives only the standardized scalar parameter,
\begin{equation}
  \widetilde\mu
  =
  \frac{\mu-\mu_R}{\sigma_R},
  \qquad
  \bm\pi(\mu)
  =
  \operatorname{softmax}
  \left(
    \mathcal R_{E2}(\widetilde\mu)
  \right)
  =
  [\pi_S,\pi_H,\pi_P]^\top .
  \label{eq:app_e2_parameter_only}
\end{equation}
The trajectory-level preference is evaluated once and remains fixed
throughout the subsequent rollout. The regime graph contains only adjacent transitions,
\begin{equation}
  \mathcal E_{\rm adj}
  =
  \bigl\{
    \{S,H\},
    \{H,P\}
  \bigr\},
  \label{eq:app_adjacency}
\end{equation}
so direct $S$--$P$ fusion is excluded.
Let
$m_j(\mathcal H_n,\mu,K)\in\{0,1\}$
denote the validation-supported applicability of specialist $j$ for the
given initialization and rollout horizon.
The regularized E2 weight restricted to applicable specialists is
\begin{equation}
  \widetilde\pi_j
  =
  \frac{
    m_j\pi_j
  }{
    \sum_{k\in\{S,H,P\}}m_k\pi_k+\varepsilon
  }.
  \label{eq:app_masked_e2}
\end{equation}
For $Z=\sum_k m_k\pi_k>0$, these regularized weights sum to $Z/(Z+\varepsilon)$, rather than exactly one. They preserve the ranking of applicable specialists; the selected pair is normalized separately below. A single applicable specialist yields the Top-1 prediction.
Top-2 fusion is considered only for an applicable pair belonging to
$\mathcal E_{\rm adj}$.
If no specialist is applicable, the hierarchy abstains from prediction. This is the general interface contract, not an experimentally validated out-of-distribution detector. In the reported centered-square overlap implementation, the evaluation cache specifies a fixed candidate pair, $(S,H)$ or $(P,H)$, and $K=24$. E2 Top-1 chooses within this supplied pair and T2-C blends its two independently generated trajectories. These experiments evaluate conditional routing and assembly on prescribed overlap support; they do not test a learned history-dependent mask, the three-applicable case, or abstention coverage on arbitrary new inputs.

We use the canonical ordered pairs $(r_1,r_2)=(S,H)$ and $(P,H)$, including in all tables reporting candidate weights. Reversing the order replaces $\alpha$ by $1-\alpha$ and reverses the corresponding logit sign. For an ordered selected pair $(r_1,r_2)$, the pair-normalized E2
preference is
\begin{equation}
  \bar\pi_1
  =
  \frac{\pi_{r_1}}
       {\pi_{r_1}+\pi_{r_2}},
  \qquad
  \bar\pi_2
  =
  1-\bar\pi_1,
  \qquad
  \ell^{E2}
  =
  \log
  \frac{\bar\pi_1+\varepsilon}
       {\bar\pi_2+\varepsilon}.
  \label{eq:app_pair_logit}
\end{equation}


\paragraph{Chart-independent physical-history descriptor.}
The centered-square implementation uses eight descriptor channels and one parameter input. These features are computed from the initial three physical states at $t_{n-2}<t_{n-1}<t_n$, with pressure mapped to the common zero-mean gauge. They are independent of the specialists' POD coordinates. All quantities below use the nondimensional data conventions of the experiment, and $\epsilon=10^{-12}$ is the numerical floor.

Define $A_\Omega=\sum_i w_i$, $U_j=\|\bm u_j\|_{M_u}^2$, $P_j=\|p_j^\circ\|_{M_p}^2$, and
\begin{align}
 h_j&=\max(t_j-t_{j-1},\epsilon),&
 s_j^u&=\sqrt{\max(U_j,\epsilon)},&
 s_j^p&=\sqrt{\max(P_j,\epsilon)},\\
 v_j&=(\bm u_j-\bm u_{j-1})/h_j,&
 g_j&=\frac{\log\max(U_j,\epsilon)-\log\max(U_{j-1},\epsilon)}{h_j}.
\end{align}
The ordered raw descriptor is
\begin{equation}
\bm s_n=\operatorname{col}\left(
\begin{array}{l}
 \tfrac12\log\max(U_n/A_\Omega,\epsilon),\\
 \tfrac12\log\max(P_n/A_\Omega,\epsilon),\\
 \log(1+\|v_n\|_{M_u}/(s_n^u+\epsilon)),\\
 \log(1+\|v_{n-1}\|_{M_u}/(s_{n-1}^u+\epsilon)),\\
 \tanh(g_n),\\
 \tanh(g_n-g_{n-1}),\\
 \log(1+\|p_n^\circ-p_{n-1}^\circ\|_{M_p}/[h_n(s_n^p+\epsilon)]),\\
 \log(1+\|v_n-v_{n-1}\|_{M_u}/(s_n^u+\epsilon))
\end{array}\right).
\label{eq:app_t2c_descriptor}
\end{equation}
The last channel measures a difference of velocity rates: it is not divided again by a time interval and is not the second-derivative descriptor in an acceleration discretization. The log and hyperbolic-tangent transformations above are part of the actual feature construction.

For the gate, the nine raw features $[\mu;\bm s_n]$ are standardized using the means and population standard deviations of training windows only. Any feature standard deviation below $10^{-8}$ is replaced by one. We denote the resulting vector by $[\widetilde\mu;\bm d_n]$; its parameter normalization is fitted for the gate and need not equal the trajectory-level E2 normalization. The architecture-matched $\mu$-only control sets all eight descriptor channels to zero after standardization. No state from inside the prediction window is used to construct these features.


\paragraph{Physical-space T2-C fusion.}
The logit-correction network receives
$[\widetilde\mu;\bm d_n]\in\mathbb R^9$ and uses two hidden layers of width 64 with SiLU
activations followed by a scalar linear output (9--64--64--1; 4,865 trainable parameters).
The final layer is initialized at zero, so the initial correction vanishes
and the fusion weight reduces to the E2 pair preference.
Specifically,
\begin{equation}
  \Delta_\psi
  =
  c_\psi
  \left(
    [\widetilde\mu;\bm d_n]
  \right),
  \qquad
  \alpha
  =
  \operatorname{sigmoid}
  \left(
    \ell^{E2}+\Delta_\psi
  \right).
  \label{eq:app_t2c_gate}
\end{equation}

The two selected specialists encode the same physical history in their
respective local charts and then evolve independently.
The scalar $\alpha$ is fixed throughout the prediction window and is shared
by velocity and pressure.
After reconstruction on the common physical mesh, T2-C forms
\begin{equation}
  \widehat{\bm q}_{\alpha}^{\,c}(t)
  =
  \alpha\,
  \widehat{\bm q}_{1}^{\,c}(t)
  +
  (1-\alpha)\,
  \widehat{\bm q}_{2}^{\,c}(t),
  \qquad
  c\in\{u,p\},
  \label{eq:app_t2c_field_fusion}
\end{equation}
where the pressure field is gauge corrected.
The fused field is an output only and is never projected back into either
specialist.

The correction network is trained directly on the reconstructed physical
fields. For rollout windows $w$ and prediction steps $k$,
\begin{equation}
  \mathcal L_{\rm T2C}
  =
  \operatorname{mean}_{w,k}
  \left[
    \frac{
      \left\|
        \widehat{\bm u}_{\alpha}^{(w,k)}
        -
        \bm u_{\rm true}^{(w,k)}
      \right\|_{M_u}^2
    }{
      \left\|
        \bm u_{\rm true}^{(w,k)}
      \right\|_{M_u}^2+\varepsilon
    }
    +
    \frac{
      \left\|
        \widehat p_{\alpha}^{\circ,(w,k)}
        -
        p_{\rm true}^{\circ,(w,k)}
      \right\|_{M_p}^2
    }{
      \left\|
        p_{\rm true}^{\circ,(w,k)}
      \right\|_{M_p}^2+\varepsilon
    }
  \right].
  \label{eq:app_t2c_loss}
\end{equation}
Optimization and checkpoint-selection details are reported in
Appendix~\ref{app:implementation}.

\paragraph{Numerical convention in the centered-square evaluator.}
The archived implementation normalizes the selected E2 probabilities, clips the first probability to $[10^{-6},1-10^{-6}]$, and then forms its logit. Its squared relative field loss uses $\max(\|q_{\rm true}\|_M^2,10^{-12})$ as denominator rather than an additive regularizer. Feature standard deviations below $10^{-8}$ are replaced by one before standardization. These are implementation-specific safeguards; the generic additive-$\varepsilon$ notation above should not be read as a claim that all numerical constants or normalization rules are shared across benchmarks.
